\section{Design}
\label{sec:design}

\subsection{Cost model}
For ranks $0..n{-}1$ in ring order with sites $\sigma(i)$, the ring time is
\[
T_{\text{ring}} = 2(n-1)\,\kappa\,\max_{e}\Big[\frac{S}{n\,b_e} + \ell_e\Big],
\]
where edge $e=(\sigma(i),\sigma(i{+}1))$ is a LAN edge if both ranks share a site. Otherwise it uses the
path $r$ selects for (src, dst, \textsc{compute}). The edge bandwidth $b_e$ is the minimum, over the
links $l$ of that path, of the fair share $\max(\epsilon, c_l - g_l)/f_l$: capacity minus background
load, divided among the controlled flows on $l$. $\ell_e$ is the path latency. $\kappa$ is a calibration
factor, initialised from a measured run and refined online from rank-0 collective timings. Every $K$
steps, each rank uploads $B/n$ bytes to site $s$ over the path $r$ selects for (site, $s$,
\textsc{storage}); while that flush is in flight, ring edges that share a link with it get a halved share.
With flush duration $D$ and $m=\min(K, D/T_{\text{ring,flush}})$ overlapped steps, the expected step time is
\[
E[T] = C + \frac{(K-m)\,T_{\text{ring}} + m\,T_{\text{ring,flush}}}{K} + \frac{\text{stall}}{K},
\]
where the stall is any part of the flush that does not finish within $K$ steps. Migrating costs a drain
(final checkpoint), a measured restart time, and a restore in which each rank pulls the state from the
old storage site. Before relying on the ring structure, we checked it with a microbenchmark
(\texttt{bench/m3\_ring\_micro.py}, \CringTrials{} trials). With the same per-site counts, interleaving
sites in rank order raises the all-reduce time from \CringContig\ to \CringInter\,s, a factor of
\CringRatio{}$\pm$\CringRatioCi. The uncalibrated model ($\kappa=1$) predicts \CringPredContig\ and
\CringPredInter\,s, a factor of \CringPredRatio. Calibration on this testbed gives a compute time of
\CcalComputeS\,s per step, \CcalGradMB\,MB of gradients, a \CcalCkptMB\,MB training state, a
\CcalRestartS\,s restart, and $\kappa_A=\CcalKappaA$, $\kappa_B=\CcalKappaB$. The two planes need
different corrections, which is why $\kappa$ is kept per plane (D-014).

\subsection{Joint planner}
With \cfg{6} ranks, \cfg{3} sites, at most \cfg{3} ranks per site and a residency-constrained $s$, the space is small. It
contains the contiguous placements and their distinct cyclic orders, \cfg{2} storage sites, and $2^k$ plane
assignments over the $k\le \cfg{5}$ route keys that a given $(\pi,s)$ actually uses. The planner therefore
enumerates the space exactly in milliseconds. On every \cfg{1}\,s controller tick it (i) re-optimises routes
for the current placement, and (ii) proposes a re-placement only if the predicted saving over a
horizon exceeds $\cfg{1.5}\times$ the predicted migration cost and the per-step gain exceeds \cfg{3}\%. Route changes
below the same \cfg{3}\% threshold are suppressed to avoid churn.

\subsection{Baselines}
\begin{itemize}
\item \textbf{static}: default placement (two ranks per site), checkpoint storage off rank~0's site, base
BGP policy (plane A, the low-latency ``private wave''); never re-plans.
\item \textbf{ecmp}: as static, but gateways run iBGP multipath and hash flows across both planes.
\item \textbf{independent}: placement chosen once by the ML team (the best placement for the nominal network,
ignoring checkpoints); storage by fixed policy; routes re-optimised continuously by a B4-style TE
MILP that minimises maximum link utilisation for the \emph{measured} per-prefix demand matrix, with a
latency tie-break. This is our stand-in for ``independently tuned subsystems''.
\item \textbf{oracle}: the joint planner fed ground-truth link state and background schedule, an upper
bound on what the model allows.
\end{itemize}

\subsection{Signals}
The controller sees: interface and per-prefix flow counters (\texttt{util}); device-reported link rate and
oper-status (\texttt{speed}, like ifSpeed via gNMI); \cfg{5}\,Hz RTT probes on every WAN link (\texttt{rtt});
rank-0 step/collective telemetry (\texttt{coll}, for online $\kappa$); and checkpoint flush statistics
(\texttt{ckpt}). The ablation removes each signal in turn, and separately disables migration.
